\section*{第 \ref{ch:b2:structure-determination} 章——结构测定：\texorpdfstring{$^{13}$C}{13C} 核磁共振、质谱与红外光谱联用}

\begin{solution}{exo:b2:structure-determination:1}
1,2-二甲苯：4 种（\ce{CH3}；C1/C2；C3/C6；C4/C5）。1,3-：5 种（\ce{CH3}；C1/C3；C2；C4/C6；C5）。
1,4-：3 种（\ce{CH3}；C1/C4；C2/C3/C5/C6）。
\end{solution}

\begin{solution}{exo:b2:structure-determination:2}
\ce{CH3CH(OH)CH2CH3}。DEPT-135：C1（\ce{CH3}）向上，C2（CH）向上，C3（\ce{CH2}）向下，C4（\ce{CH3}）向上。
DEPT-90：只有 C2。
\end{solution}

\begin{solution}{exo:b2:structure-determination:3}
209：酮羰基；170：酯（或羧酸）羰基；129：芳香碳；64：与
氧相连的碳；14：烷基 \ce{CH3}。
\end{solution}

\begin{solution}{exo:b2:structure-determination:4}
$(2 \times 8 + 2 - 8)/2 = 5$：一个苯环（4）和一个 \ce{C=O}（1）。例如苯甲酸甲酯，
\ce{C6H5COOCH3}，或 4-甲基苯甲酸。
\end{solution}

\begin{solution}{exo:b2:structure-determination:5}
一个不饱和度，属酮（\qty{1715}{cm^{-1}}，209 处谱线，\omterm{def:b2:structure-determination:dept}{季碳}）；一个 \ce{CH2} 和两个
\ce{CH3}，四种不同的碳：丁-2-酮，\ce{CH3COCH2CH3}。
\end{solution}

\begin{solution}{exo:b2:structure-determination:6}
碳谱：戊-2-酮给出 5 条谱线，对称的戊-3-酮给出 3 条。质子谱：戊-2-酮有一个
三质子单峰（\ce{CH3CO}）；戊-3-酮只有一个四重峰和一个三重峰。质谱：戊-2-酮
\omterm{def:b2:mass-spec-atomic:spectrum}{基峰}在 43（\ce{CH3CO+}），\omterm{def:b2:mass-spec-atomic:fragmentation}{麦氏重排}离子在 58；戊-3-酮的\omterm{def:b2:mass-spec-atomic:spectrum}{基峰}在 57
（\ce{C2H5CO+}）。
\end{solution}

\begin{solution}{exo:b2:structure-determination:7}
\ce{C4H8O2}，88：四重峰–三重峰组合为 \ce{OCH2CH3}（其中 \ce{CH2} 位于 4.12，与氧相连），
单峰为 \ce{CH3CO}。乙酸乙酯，\ce{CH3COOCH2CH3}。
\end{solution}

\begin{solution}{exo:b2:structure-determination:8}
1,2-二氯苯：三条谱线（C1/C2，C3/C6，C4/C5）。1,4-二氯苯：两条（C1/C4，四个
CH）。DEPT-90 中，前者给出两条 CH 谱线，后者给出一条。
\end{solution}

\begin{solution}{exo:b2:structure-determination:9}
C2 为手性碳原子。将 C3 上两个质子中的任一个换成氘，得到
两个非对映异构体，而不是对映异构体：这两个质子是非对映异位的，处于不同环境中。
它们的化学位移可以不同，既彼此耦合，也与邻近质子耦合，因此
\ce{CH2} 给出复杂的多重峰，而不是简单的五重峰。
\end{solution}

\begin{solution}{exo:b2:structure-determination:10}
五个不饱和度；一个单取代苯环（三条 CH 谱线和一条 C 谱线，位于 137.0；五个芳香
质子）；一个酮（200.8，C）；一个乙基（\ce{CH2}：31.8，质子 2.98 处四重峰，与羰基相邻；\ce{CH3}
：8.3）。1-苯基丙-1-酮，\ce{C6H5COCH2CH3}。与苯环共轭使 \ce{C=O} 的
$\pi$ 电子离域，削弱了该键，吸收带低于饱和酮的 \qty{1715}{cm^{-1}}。
% ledger: ir:CO-ketone
\end{solution}

\begin{solution}{exo:b2:structure-determination:11}
戊酸：在 2500 至 \qty{3300}{cm^{-1}} 间有很宽的 \ce{O-H} 吸收带，另有一个质子约在 11--12
ppm 处。丁酸甲酯：三质子单峰位于与氧相连碳上质子的范围内
（3.3--4.5）。丙酸乙酯：两个四重峰，分别与氧和羰基相邻。乙酸丙
酯：与羰基相邻的三质子单峰（2.0--2.4），与氧相连的三重峰。
% ledger: ir:OH-acid, nmrr:acid, nmrr:alcohol, nmrr:ketone
\end{solution}

\begin{solution}{exo:b2:structure-determination:12}
两条谱线在 \omterm{def:b2:structure-determination:dept}{DEPT} 中均消失，故均为\omterm{def:b2:structure-determination:dept}{季碳}，仅凭 \omterm{def:b2:structure-determination:dept}{DEPT} 无法区分二者。化学位移
则可以：酯羰基位于 167--171 附近，芳香碳位于 128 至 137 之间（本章
图表）。若环碳位于 166.5、羰基位于 130.6，二者都将远超出各自的范围；
正确的归属恰好相反。
\end{solution}

\begin{solution}{pb:b2:structure-determination:1}
\textbf{1.} $3.1/31.7 \approx 9.8~\%$；$9.8/1.11 \approx 9$ 个碳原子。
\textbf{2.} 质量为偶数：不含氮（或含两个氮）。
\textbf{3.} $150 - 108 = 42$：\ce{H10O2} 符合（10 + 32）：\ce{C9H10O2}。
\textbf{4.} 若为十个碳，则 $M+1 \approx 11~\%$ 的 $M$，而非 9.8~\%；且红外谱显示
存在羰基，它需要另一个氧原子，此外还有 \ce{C-O} 带。
\textbf{5.} $(2 \times 9 + 2 - 10)/2 = 5$.
\textbf{6.} $108 + 10 \times 1.007825 + 2 \times 15.994915 \approx 150.0681$.
\textbf{7.} \ce{C=O} 伸缩振动位于饱和酯的位置（约 \qty{1735}{cm^{-1}}）。
\textbf{8.} 酯的 \ce{C-O} 伸缩振动。
\textbf{9.} 芳香（或烯烃）\ce{C-H} 键。
\textbf{10.} 含 \ce{O-H} 基团：不是醇，也不是羧酸。
\textbf{11.} 否：与环共轭的羰基吸收波数更低；此处吸收位于
饱和酯的位置，因此有一个不共轭的基团将其与环隔开。
\textbf{12.} 170.70：酯 \ce{C=O}；136.14：带取代基的环碳；128.56 和 128.24：
芳香 CH；66.24：与氧相连的 \ce{CH2}；20.82：\ce{CH3}。
\textbf{13.} 66.24，即 \ce{CH2}。
\textbf{14.} 170.70 和 136.14，即两个不含氢的碳。
\textbf{15.} 128.56 和 128.24。
\textbf{16.} 四个：带取代基的碳、两个邻位、两个间位和对位 CH。预期有三
种 CH；其中两种化学位移过于接近，在此场强下无法分辨，共用一条谱线。
\textbf{17.} 未必：峰高不等于计数（见\cref{prop:b2:structure-determination:intensities}），
不过此处这条谱线很可能包含两种重叠的 CH。
\textbf{18.} 7.33：五个芳香质子；5.09：\ce{OCH2}；2.06：\ce{CH3CO}。
\textbf{19.} 它们的相邻原子（对 \ce{CH2} 而言为环碳和氧，对 \ce{CH3} 而言为羰基
碳）均不含氢。
\textbf{20.} 它与酯中电负性大的氧及苯环相连，两者都使其去屏蔽。
\textbf{21.} \ce{C6H5CH2OCOCH3}，乙酸苄酯。
\textbf{22.} 在该异构体中，\ce{CH2} 不与氧相连，故不可能出现在 5.09，而
与氧相连的甲基应出现在 3.3--4.5 范围内，而非与羰基相邻的 2.06 处。
% ledger: nmrr:alcohol, nmrr:ketone
\textbf{23.} 108：经重排失去 42（烯酮，\ce{CH2=C=O}）得到 \ce{C7H8O}
的自由基正离子；91：\ce{C7H7+}，苄基正离子；43：\ce{CH3CO+}。
\textbf{24.} 共七种碳（\ce{C=O}、取代的环碳、邻位、间位和对位 CH、
\ce{CH2}、\ce{CH3}）：\textbf{七条谱线，其中一条（\ce{CH2}）向下}，出现在 DEPT-135 中。
\end{solution}
